% Generated draft: requires mathematical and source review.
\documentclass[11pt]{article}
\usepackage[margin=25mm]{geometry}
\usepackage{fontspec}
\setmainfont{Latin Modern Roman}
\usepackage{amsmath,amssymb,mathtools,unicode-math}
\setmathfont{Latin Modern Math}
\usepackage{enumitem,needspace,xurl,hyperref}
\setlength{\parindent}{0pt}
\setlength{\parskip}{5pt}
\setlength{\emergencystretch}{3em}
\newcommand{\N}{\mathbb{N}}
\newcommand{\Z}{\mathbb{Z}}
\newcommand{\Q}{\mathbb{Q}}
\newcommand{\R}{\mathbb{R}}
\newcommand{\C}{\mathbb{C}}
\newcommand{\F}{\mathbb{F}}
\newcommand{\A}{\mathbb{A}}
\newcommand{\PP}{\mathbb{P}}
\newcommand{\E}{\mathbb{E}}
\newcommand{\eps}{\varepsilon}
\newcommand{\dd}{\,\mathrm{d}}
\begin{document}
\section*{UnsolvedMath ID: 800163}
\textbf{Source problem code:} catalogue entry
\subsection{Definitions and mathematical statement}
Let \(\zeta(s)\) denote the meromorphic continuation of the Riemann zeta function, initially defined by \(\zeta(s)=\sum_{n=1}^{\infty}n^{-s}\) for \(\operatorname{Re}(s)>1\). The supplied record does not define \(I_k(T)\). Under the conventional unweighted interpretation, put
\[
I_k(T):=\int_0^T\left|\zeta\!\left(\tfrac12+it\right)\right|^{2k}\,\dd t,
\qquad k\in\Z_{\geq1},\quad T\geq2.
\]
This normalization and these integration limits require confirmation against the source.

The problem attributed in the supplied excerpt to Nathan Ng (Lethbridge) asks whether the conjectural method called the ``recipe'' in reference [CFKRS] can be made rigorous for the fourth moment \(I_2(T)\). Under the interpretation above, the requested result is a proof following that recipe of an asymptotic formula
\[
I_2(T)=T P_4(\log T)+O\!\left(T^{1-\delta}\right)
\qquad (T\to\infty),
\]
for some constant \(\delta>0\), where \(P_4\in\R[X]\) is the fixed degree-four polynomial giving the full fourth-moment main term predicted by the recipe. The implied constant must be independent of \(T\). All coefficients of the main-term polynomial must be accounted for, and the error must save a positive power of \(T\).

The record presents this as a request for a new proof of an existing asymptotic formula. It does not supply the coefficients of \(P_4\), a definition of the recipe, or the full bibliographic identification of [CFKRS]. Consequently, the displayed formulation is conditional on confirming the normalization of \(I_k(T)\), and the precise main term and required methodological steps remain to be checked against the cited source.
\subsection{Short English statement}
Make the [CFKRS] recipe rigorous for the fourth moment of the Riemann zeta function, obtaining its complete polynomial main term with an error saving a positive power of \(T\). The supplied excerpt leaves the moment normalization, polynomial coefficients, and precise recipe unspecified.
\subsection{Sources}
\begin{itemize}
\item Local UnsolvedMath registry, permanent ID \(800163\), source code \texttt{catalogue entry}. The supplied record labels the problem ``open''; this is a supplied registry classification, not an independently verified status assessment.
\item Hamieh, Kadiri, Ng, \emph{Moments of L-functions Problem List} (2024), Problem 13, in the portion concerning moments of the Riemann zeta function. Supplied source URL: \url{https://arxiv.org/abs/2405.17800}. The background records access on 2026-10-16 and extraction by \texttt{arxiv\_latex}; these are supplied provenance details, not a claim of independent source access or verification.
\item The problem excerpt explicitly attributes the question to Nathan Ng (Lethbridge). The registry lists Hamieh, Kadiri, Ng as proposers and gives 2024 as the proposed year.
\item Reference [CFKRS] is cited in the supplied statement as introducing the conjectural ``recipe.'' Its full bibliographic entry and a direct reference for the existing fourth-moment asymptotic formula are not supplied. No additional bibliographic identification has been inferred.
\end{itemize}
\end{document}
